% !TeX root = ../main.tex

\chapter{Centered Blocks}

\section{Changing the Exponential Base}
% =================================================

The choice of base \(b\) can be simplified.

Since

\[
b=e^\lambda,
\]

we have

\[
b^{x^n}
=
e^{\lambda x^n}.
\]

Therefore it is enough to study

\[
\boxed{
\phi_n(x)=e^{\lambda x^n}.
}
\]

The parameter

\[
\lambda=\ln b
\]

contains all information about the base.

There are three fundamentally different cases:

\[
\lambda>0
\quad\Longleftrightarrow\quad
b>1,
\]

\[
\lambda<0
\quad\Longleftrightarrow\quad
0<b<1,
\]

and

\[
\lambda=0
\quad\Longleftrightarrow\quad
b=1.
\]

The case \(b=1\) is trivial because

\[
1^{x^n}=1
\]

for every \(n\).

We shall therefore exclude it.

% =================================================
\section{Taylor Expansion of a Block}
% =================================================

The representation

\[
\phi_n(x)=e^{\lambda x^n}
\]

reveals another important structure.

Recall that

\[
e^y
=
\sum_{k=0}^{\infty}\frac{y^k}{k!}.
\]

Substituting

\[
y=\lambda x^n,
\]

we obtain

\[
e^{\lambda x^n}
=
\sum_{k=0}^{\infty}
\frac{\lambda^k x^{nk}}{k!}.
\]

Therefore

\[
\boxed{
\phi_n(x)
=
1
+
\lambda x^n
+
\frac{\lambda^2}{2!}x^{2n}
+
\frac{\lambda^3}{3!}x^{3n}
+\cdots.
}
\]

Notice something important:

\[
\phi_n
\]

contains only powers whose exponents are multiples of \(n\):

\[
n,2n,3n,4n,\ldots.
\]

Thus divisibility appears once again.

% =================================================
\section{Taylor Coefficients and Divisors}
% =================================================

Consider formally

\[
F(x)
=
\sum_{n\geq1}c_ne^{\lambda x^n}.
\]

Using the Taylor expansion,

\[
F(x)
=
\sum_{n\geq1}
c_n
\sum_{k\geq0}
\frac{\lambda^k}{k!}x^{nk}.
\]

Suppose we want the coefficient of \(x^m\).

A contribution appears exactly when

\[
nk=m.
\]

Therefore

\[
n\mid m.
\]

For \(m\geq1\), the coefficient of \(x^m\) is consequently

\[
\boxed{
[x^m]F(x)
=
\sum_{n\mid m}
c_n
\frac{\lambda^{m/n}}{(m/n)!}.
}
\]

This is another divisor sum.

Thus divisibility enters the theory in two completely different ways:

\[
\text{composition}
\quad\text{and}\quad
\text{Taylor expansion}.
\]

This repetition suggests that the divisor structure is not merely
accidental.

% =================================================
\section{Centered Blocks}
% =================================================

There is one inconvenience with the blocks

\[
\phi_n(x)=e^{\lambda x^n}.
\]

At \(x=0\),

\[
\phi_n(0)=1
\]

for every \(n\).

It is therefore useful to define the centered blocks

\[
\boxed{
\psi_n(x)
=
e^{\lambda x^n}-1.
}
\]

Now

\[
\psi_n(0)=0.
\]

Their Taylor expansion is

\[
\psi_n(x)
=
\lambda x^n
+
\frac{\lambda^2}{2!}x^{2n}
+
\frac{\lambda^3}{3!}x^{3n}
+\cdots.
\]

The lowest-degree term is exactly

\[
\boxed{
\lambda x^n.
}
\]

This gives the family a triangular relationship with the ordinary
monomials.

% =================================================
